% Part: first-order-logic
% Chapter: sequent-calculus
% Section: rules-and-proofs
%READERNOTE{SOL6-A144: प्रमेयाच्या क्रमवर्तीत फक्त निर्दिष्ट सत्यतामूल्यांच्या जागी सूत्र असते; इतर जागा रिकाम्या असतात. हा आवश्यक scope स्पष्ट केला. व्युत्पादन सांत वृक्ष आहे आणि त्यात तार्किक व संरचनात्मक नियम दोन्ही येतात. दाखवलेले Gamma उदाहरण सांत क्रमिकेसाठी; अनंत आधारसंचासाठी आधीच्या सांत-उपसंच व्याख्येचा वापर होतो.}

\documentclass[../../../include/open-logic-section]{subfiles}

\begin{document}

\olfileid{mvl}{seq}{rul}

\olsection{नियम आणि \usetoken{P}{derivation}}

पुढे $\Gamma, \Delta, \Pi, \Lambda$ या !!{sentence}s च्या
सांत क्रमिका दर्शवतात असे समजू.

\begin{defn}[क्रमवर्ती]
\emph{$n$-बाजूंची क्रमवर्ती} म्हणजे पुढील रूपाची अभिव्यक्ती:
\[
\Gamma_1 \nSequent \dots \nSequent \Gamma_n
\]
येथे प्रत्येक $\Gamma_i$ ही भाषा~$\Lang L$ मधील
!!{sentence}s ची सांत, शक्यतो रिकामी, क्रमिका असते.
\end{defn}

\begin{defn}[आरंभीची क्रमवर्ती]
\emph{$n$-बाजूंची आरंभीची क्रमवर्ती} म्हणजे भाषेतील कोणत्याही
!!{sentence}~$!A$ साठी $!A \nSequent \dots \nSequent !A$
या रूपाची $n$-बाजूंची क्रमवर्ती.

भाषेत $0$-स्थानी संयोजक~$\star$, म्हणजे विधानीय स्थिरांक,
असेल तर $\dots \nSequent \star \nSequent \dots$ ही क्रमवर्तीही
आरंभीची मानतो. यात $\tf{\star} \in V$ शी संबंधित
सत्यतामूल्याच्या जागी $\star$ असते; इतर जागा रिकाम्या असतात.
\end{defn}

$n$-मूल्य तर्कशास्त्र~$\Log{L}$ मधील प्रत्येक संयोजकासाठी,
त्या संयोजकाला $\Log{L}$ मध्ये मिळू शकणाऱ्या प्रत्येक
सत्यतामूल्यासाठी एक तार्किक नियम असतो. $\Log{L}$ च्या
$n$-बाजूंच्या क्रमवर्ती कलनातील !!^{derivation}s ही क्रमवर्तींची
सांत वृक्षरूप मांडणी असते: सर्वांत वरच्या क्रमवर्ती आरंभीच्या असतात;
आणि एखादी क्रमवर्ती एक किंवा अधिक क्रमवर्तींच्या खाली असेल,
तर ती या कलनातील निगमन नियमाने त्यांच्यापासून योग्यरीत्या
निष्पन्न झालेली असावी लागते. यात संयोजकांचे तार्किक नियम
आणि पुढील विभागात दिलेले संरचनात्मक नियम दोन्ही येतात.

\begin{defn}[प्रमेये]
!!^a{sentence}~$!A$ हे $n$-मूल्य तर्कशास्त्र~$\Log{L}$ चे
\emph{प्रमेय} तेव्हा असते, जेव्हा $\Log{L}$ मधील प्रत्येक निर्दिष्ट
सत्यतामूल्याशी संबंधित जागेत $!A$ असलेल्या $n$-बाजूंच्या
क्रमवर्तीचे !!a{derivation} असते; या क्रमवर्तीतील
इतर सर्व जागा रिकाम्या असाव्यात. $!A$ प्रमेय असेल तर
$\Proves[\Log{L}] !A$ आणि नसेल तर $\Proves/[\Log{L}] !A$ लिहितो.
\end{defn}

\begin{defn}[!!^{derivability}]
$n$-मूल्य तर्कशास्त्र~$\Log{L}$ मध्ये !!^a{sentence}~$!A$
हे !!{sentence}s च्या संच~$\Gamma$ पासून \emph{!!{derivable}}
असते, $\Gamma \Proves[\Log{L}] !A$, तेव्हा आणि केवळ तेव्हाच,
जेव्हा एखादा ससीम उपसंच~$\Gamma_0 \subseteq \Gamma$
आणि $\Gamma_0$ मधील !!{sentence}s ची एखादी क्रमिका~$\Gamma_0'$
अशी असते की पुढील क्रमवर्तीचे !!a{derivation} असते:
\[ \Lambda_1 \nSequent \dots \nSequent \Lambda_n \]
येथे जागा~$i$ निर्दिष्ट सत्यतामूल्याशी संबंधित असेल तर
$\Lambda_i$ हे $!A$ असते; अन्यथा ते $\Gamma_0'$ असते.
$!A$ हे $\Gamma$ पासून !!{derivable} नसेल तर
$\Gamma \Proves/ !A$ असे लिहितो.
\end{defn}

उदाहरणार्थ, $3$-मूल्य \L ukasiewicz तर्कशास्त्रासाठी
$3$-बाजूंचे क्रमवर्ती कलन आहे. $3$-बाजूंच्या क्रमवर्ती
$\Gamma \nSequent \Pi \nSequent \Delta$ मध्ये $\Gamma$
हे~$\False$ शी, $\Delta$ हे~$\True$ शी आणि $\Pi$
हे~$\Undef$ शी संबंधित असतात. स्वयंसिद्धके
$!A \nSequent !A \nSequent !A$ ही आहेत. फक्त $\True$
निर्दिष्ट असल्यामुळे, $\Gamma$ ही सांत क्रमिका असल्यास
$\Gamma \Proves[\LogLuk[3]] !A$
तेव्हा आणि केवळ तेव्हाच, जेव्हा $\Gamma \nSequent \Gamma \nSequent !A$
या क्रमवर्तीचे !!a{derivation} असते. ($\Undef$ देखील निर्दिष्ट
असते, तर $\Gamma \nSequent !A \nSequent !A$ या क्रमवर्तीचे
!!a{derivation} आवश्यक झाले असते.)
अनंत आधारसूत्रसंचासाठी, आधीच्या व्याख्येनुसार त्यातील सांत
उपसंचाच्या क्रमिकेचा वापर करायचा; अनंत क्रमवर्तीचा नाही.

\end{document}
